% ============================================================
% PREFACE
% front_matter/preface.tex
% ============================================================

\chapter*{Preface}
\addcontentsline{toc}{chapter}{Preface}
\markboth{Preface}{Preface}

\chapterquote{A person who never made a mistake never tried
anything new.}{Albert Einstein}

\dropcap{T}here was a time when the greatest scientific minds 
in human history --- including \newline{}esla lying awake at night 
visualising entire machines in his mind before ever touching 
a tool, Newton sitting beneath an apple tree, Faraday wrapping 
wire around an iron ring in his basement laboratory --- did not 
study Physics to pass an examination. They studied Physics 
because they could not help themselves. The universe presented 
them with questions they were physically unable to leave 
unanswered.

\vspace{0.5cm}

\textbf{Why this book was written}

\vspace{0.3cm}

\textit{The Physics Codex} was written because the students I
teach --- brilliant, capable, curious young Nigerians preparing
for UTME and entering university --- deserve better than what
most textbooks give them. They deserve explanations that make
them understand, not just recall. They deserve to feel the
satisfaction of a concept clicking into place. They deserve to
read about Newton's second law and think: \textit{yes, I see
it --- I have felt that my whole life and now I know what it
is called and why.}

This book is built on one principle: \textbf{understanding
first, examination second.} Every concept in these pages is
explained from its foundations. Nothing is stated without
reason. No formula appears without derivation or intuition.
Every topic is connected to the physical world you can see,
touch, and experience --- because Physics is not abstract
mathematics; it is the language the universe uses to describe
itself.

\vspace{0.5cm}

\textbf{Who this book is for}

\vspace{0.3cm}

This textbook is written for two groups of students, and
deliberately bridges the gap between them:

\begin{itemize}
  \item \textbf{Senior Secondary School students (SS2 and SS3)}
  preparing for UTME who want to understand Physics deeply
  enough that examination questions feel straightforward ---
  because when you truly understand a concept, no question
  about it can surprise you.

  \item \textbf{First-year university students} who need a
  rigorous, complete foundation in Physics at a level that
  connects their secondary school knowledge to the demands of
  university coursework.
\end{itemize}

The language is accessible enough for an SS2 student. The
depth is sufficient for a first-year university student. That
is not a contradiction --- it is the mark of genuinely good
teaching.

\vspace{0.5cm}

\textbf{What makes this book different}

\vspace{0.3cm}

\begin{enumerate}
  \item \textbf{Every explanation begins with intuition.}
  Before any formula is introduced, the physical idea behind
  it is established using examples from everyday life ---
  including examples rooted in the Nigerian context you
  already know.

  \item \textbf{Every formula is derived, not just stated.}
  A formula you understand is a formula you will never forget.
  A formula you merely memorise will betray you under pressure.

  \item \textbf{Every diagram is precise and instructive.}
  The diagrams in this book are not decorative. Each one is
  designed to show something that words alone cannot.

  \item \textbf{Every question is original.}
  The practice questions and worked examples in this book were
  created specifically for it. They are not borrowed from
  other textbooks or past examination papers. They test
  understanding, not familiarity.

  \item \textbf{Historical context is provided.}
  You will learn not just what the laws of Physics are, but
  who discovered them, why they were looking, what problem
  they were trying to solve. Science is a human story, and
  understanding that story makes the science itself more
  memorable and more meaningful.
\end{enumerate}

\vspace{0.5cm}

\textbf{A note to the student}

\vspace{0.3cm}

Read slowly. Do not rush to the formulas. When you encounter
a concept you do not immediately understand, do not skip it ---
sit with it. Re-read the paragraph. Look at the diagram.
Try the worked example before reading the solution. Ask why.

The great physicists were not great because they were
unusually intelligent. They were great because they refused
to accept that they did not understand something. That refusal
--- that insistence on genuine comprehension --- is available
to every single one of you.

This book is your tool. How well it serves you depends
entirely on how seriously you use it.

\vspace{1cm}

\begin{flushright}
\textit{Korede Samuel Omotosho (Falcon)}\\
\textit{Lagos, Nigeria}\\
\textit{2026}
\end{flushright}
